\section{Materials and Methods}
\label{sec:methods}
\subsection{Study Design}
The study is a supervised binary classification experiment that compares two fixed models, TF--IDF with logistic regression (LR) and fine-tuned \texttt{roberta-base}, on identical held-out reviews. The primary outcome is test macro-F1, the unweighted mean of the F1 scores of the deceptive and truthful classes. Secondary outcomes are deceptive-class precision, recall and F1, accuracy, average precision (AP), the area under the receiver operating characteristic curve (ROC-AUC), robustness to formatting changes and computational cost. Model settings, seeds, the decision threshold and the statistical tests were fixed in the Assignment~2 protocol before any model was trained; deviations from that protocol are listed in Section~\ref{sec:deviations}.

\subsection{Data Source, Labels and Provenance}
The experiments use the Deceptive Opinion Spam corpus v1.4 (OpSpam) \cite{ref22,ott2013}. It contains 1,600 English reviews of 20 Chicago hotels: 400 truthful positive reviews from TripAdvisor, 400 deceptive positive reviews written by Mechanical Turk workers, 400 truthful negative reviews from six review sites and 400 deceptive negative reviews from Mechanical Turk, with 20 reviews per hotel in each of the four cells. The positive truthful reviews are 5-star TripAdvisor reviews of at least 150 characters \cite{ref22}. Deceptive reviews were written by Mechanical Turk workers, each allowed one task, who were asked to imagine working in a hotel's marketing department and to write a fake positive review of that hotel or, for the negative set, a fake negative review of a competitor's hotel \cite{ref22,ott2013}. The target is $y=1$ for a deceptive review and $y=0$ for a truthful one. Deceptive labels are known by construction; truthful labels rest on the original authors' collection and filtering of real reviews and are therefore not independently verified.

The author-hosted archive returned HTTP 403 and the ACL supplementary file contains an older release without raw truthful texts, so the corpus was obtained from a public university GitLab mirror that includes the original README and CC BY-NC-SA 3.0 licence \cite{opspam_mirror}. The downloaded archive (SHA-256 \texttt{f6b9ce96\ldots{}a1c287}) and every review file were hashed and recorded in a data manifest. All 400 positive deceptive texts match the official ACL raw texts after whitespace normalisation, which partially authenticates the mirror. The corpus has no reviewer identifiers, star ratings or timestamps.

The original protocol named the text-containing YelpChi collection (67,395 reviews with Yelp filter labels \cite{ref21,rayana2015}) as the primary dataset. Its maintainer distributes it only on request, and no copy was obtained; the graph-feature file of the same name lacks review text and was deliberately not used as a substitute.

\subsection{Cleaning, Duplicate Audit and Partitioning}
No review was empty or had a label that disagreed with its filename prefix, so all 1,600 were retained (Table~\ref{tab:data}). For duplicate detection only, texts were normalised with Unicode NFKC and whitespace collapsing; models received the original text. Exact duplicates were grouped, and all pairs of reviews were screened for near-duplicates using lowercased word 5-shingles with a Jaccard similarity of at least 0.80.

Grouping is needed because a random review-level split would place reviews of the same hotel in training and test data. Reviews of one hotel share its name, location, amenities and staff, and both classes mention them, so a model could learn hotel-specific vocabulary that is uninformative about deception elsewhere, and its test score would overstate performance on unseen hotels \cite{kapoor2023}. Holding out whole hotels makes the test closer to the practical task of screening reviews of a property the model has not seen. Because reviewer identifiers are absent, hotels were used as the grouping unit. Reviews were joined into connected components if they shared a hotel or a (near-)duplicate link, and components were never split. Each component was keyed by its smallest review path, and components were ordered by the SHA-256 digest of \texttt{"<seed>|<key>"}. The ordered components were assigned greedily to training until 70\% of reviews were allocated, then to validation until 85\%, and the remainder to test. The procedure fails explicitly if any partition is empty or single-class. The primary partition uses seed 42; seeds 123 and 2026 generate two further partitions for sensitivity analysis.

To estimate uncertainty with more than three test hotels, an exploratory hotel-grouped 5-fold cross-validation was added: components were ordered by the seed-42 hash and dealt into five folds of four hotels, and within each fold the first three remaining hotels formed the validation set used only for RoBERTa checkpoint selection. Every review thus receives exactly one out-of-fold prediction from each model.

\subsection{Models}
\emph{Model 1 (baseline).} A scikit-learn \texttt{TfidfVectorizer} with lowercasing, word unigrams and bigrams, \texttt{min\_df}=3, \texttt{max\_df}=0.95, sublinear term frequency, at most 100,000 features and L2 normalisation, fitted on training reviews only. Its output feeds \texttt{LogisticRegression} with the liblinear solver, $C=1.0$, balanced class weights and at most 2,000 iterations. LR has no tuned hyperparameters, so the validation set is only scored for reference.

\emph{Model 2 (challenger).} \texttt{FacebookAI/roberta-base} \cite{liu2019roberta} with a two-class sequence-classification head, loaded through Hugging Face Transformers \cite{wolf2020}. Inputs are truncated to 256 tokens with dynamic padding. Training uses AdamW (learning rate $2\times10^{-5}$, weight decay 0.01, excluding biases and layer-normalisation weights), batch size 16, three epochs, linear decay after a 10\% warm-up, gradient clipping at 1.0 and class-weighted cross-entropy. The checkpoint with the highest validation macro-F1 is kept, with ties resolved in favour of the earlier epoch; test data are not used for selection. Python, NumPy and PyTorch seeds are set to the partition seed and deterministic algorithms are requested; PyTorch reported one non-deterministic CUDA operation (the memory-efficient attention backward pass), so repeated RoBERTa runs may differ slightly. The 256-token limit truncates 21.2\% of clean test reviews.

Both models predict the deceptive class when its probability is at least 0.5. No resampling is applied, and no threshold is tuned. These settings were fixed before training rather than tuned, for two reasons. Tuning several hyperparameters on a 240-review validation set would itself be noisy and would multiply the number of comparisons that could be reported, and the study's question concerns default, reproducible configurations of each model family rather than the best score attainable with unlimited search. The price is that neither model is necessarily at its optimum, which is acknowledged in Section~\ref{sec:conclusion}.

\subsection{Metrics and Statistical Inference}
Precision, recall and F1 use zero-division values of zero. AP is computed from continuous scores with \texttt{average\_precision\_score} and is not a trapezoidal area. Uncertainty is estimated with a 95\% percentile bootstrap that resamples whole test hotels with replacement (2,000 replicates, NumPy seed 42); the same resampled hotels are applied to both models, which yields a paired interval for $\Delta$macro-F1 (RoBERTa minus LR).

The confirmatory comparison is the two-sided paired approximate randomisation test specified in Assignment~2: within each test hotel, the two models' prediction vectors are swapped with probability 0.5, 10,000 times (seed 42), and $p=(1+\#\{|\Delta_b|\ge|\Delta_\mathrm{obs}|\})/10{,}001$ with $\alpha=0.05$. With only three test hotels, this test has $2^3=8$ distinct swap patterns, so $p$ cannot fall much below 0.25 regardless of the effect size. Two further tests are therefore reported as exploratory: the same test with review-level swapping on the seed-42 test set, which treats reviews as exchangeable \cite{dror2018}, and the hotel-level test on the pooled cross-validation predictions, which has 20 hotels and $2^{20}$ swap patterns.

Macro-F1 was chosen as the primary outcome because it weights both classes equally and cannot be raised by favouring the majority class, which keeps the protocol meaningful for the imbalanced YelpChi data for which it was designed. AP and ROC-AUC are threshold-free and show how well each model ranks reviews; a model can rank well yet perform poorly at a fixed threshold if its scores are poorly calibrated. All bootstrap intervals are conditional on the observed test hotels and describe sampling variability only; they cannot correct errors in the corpus labels. Variation across complete reruns is assessed separately by repeating the whole pipeline with seeds 123 and 2026, which changes both the held-out hotels and the training randomness. Test predictions from different seeds are not pooled, because they overlap in hotels and would not form an independent sample.

\subsection{Robustness and Computational Cost}
\label{sec:costmethod}
Three formatting perturbations were applied to the seed-42 test reviews after training: conversion to uppercase, replacement of every whitespace run with a mixed sequence of spaces, a newline and a tab, and both together. The change in macro-F1 is reported with a paired hotel-bootstrap interval. Results are also broken down by held-out hotel and by review polarity. Cost measures are training time, peak GPU memory for RoBERTa and inference throughput: after a 100-review warm-up, the 240 test reviews are classified five times in batches of 32, and the median and range in reviews per second are reported. Throughput includes tokenisation or vectorisation but excludes disk input; LR runs on the CPU and RoBERTa on the GPU, as each would be deployed.

\subsection{Implementation and Outputs}
All experiments were run on one laptop with an NVIDIA GeForce RTX 4060 Laptop GPU (8~GB) under Windows~11, using \ENVVERSIONS. These are the versions actually used; they differ from the target versions listed in Assignment~2 (Python 3.11 and older libraries), which were not available for the installed Python release. The pipeline writes a data manifest with the hash of every review file, the partition and fold assignment of every review, every model probability on the validation and test sets, and a JSON record of software versions, timings and the selected RoBERTa epochs. A separate analysis script recomputes all metrics, intervals, tests and figures from these saved files, so every number in Section~\ref{sec:results} can be checked without retraining.

\subsection{Ethical Considerations}
The corpus is public, contains no reviewer identifiers and is used under its non-commercial licence; no personal data were collected. A false positive from a deployed detector amounts to an accusation that a genuine customer lied, so false-positive counts are reported alongside every summary score, and errors are described as disagreements with corpus labels rather than as verified judgements about individuals.

\subsection{Deviations from the Original Protocol}
\label{sec:deviations}
(1) YelpChi could not be obtained, so OpSpam, planned as an external test corpus, became the development and test corpus. (2) Hotel components replace reviewer components because OpSpam has no reviewer identifiers; author-disjointness and chronological splitting cannot be established. (3) The star-rating ablation (original secondary question) was dropped because OpSpam has no ratings; polarity folders were not converted into artificial ratings. The research questions were revised accordingly: RQ2 and RQ3 now cover robustness and cost, which were secondary outcomes in the protocol. (4) The grouped cross-validation and the review-level randomisation test were added after the protocol and are exploratory. (5) The two-annotator label audit was planned for YelpChi and was not carried out.
